\chapter{Conclusion and Future Work}

\section{Conclusion}
\ \ \ \ \, This thesis asks whether server-side adaptive optimization can act
as an implicit, CSI-free filter for an $\alpha$-stable OTA-FL channel.
The main supported claim is empirical: under the three-seed CIFAR-10
tail-index ablation with $\gamma=0.05$ and $\alpha\in\{1.1,1.3,1.5,
1.7,1.9\}$, AdaGrad-OTA (49.52--51.75\%) and Adam-OTA (42.88--44.82\%)
outperform FedAvg-OTA (~11.9\%) and FedAvgM-OTA (20.5--20.7\%) by a
large margin, which supports the hypothesis that a fractional
second-moment accumulator contracts the effective server learning rate
precisely on impulsive rounds.

The claim is scoped rather than absolute. The observed $\alpha$-trend
is non-monotonic at a fixed server learning rate; MAC at the default
$\alpha=1.3$, $\gamma=0.05$, $k=3$ operating point is empirically
neutral; and in the $\alpha=2$ CIFAR-10 comparison FedAvgM-OTA
remains the strongest method (64.42\%), which shows that fractional
second-moment normalization is not costless in the Gaussian regime.
The main output is therefore a reproducible PyTorch simulator with
$\alpha$-stable tail-index experiments, MAC robustness comparisons,
and a one-dimensional intuition for the shrinkage mechanism, rather
than a complete theoretical proof or line-by-line reproduction of the
reference paper.

\section{Limitations}
\ \ \ \ \, The project remains narrower than the full reference paper in four
ways. First, it provides a one-dimensional closed-form intuition
rather than a full non-convex convergence proof. Second, the wireless
channel is a simulation oracle and does not model mobility, multi-cell
interference, synchronization errors, CSI acquisition, or power
control. Third, the project does not include a software-defined-radio
or real wireless testbed validation. Fourth, the empirical evidence is
bounded by three seeds and a single global server learning rate per
experiment, so the non-monotonic $\alpha$-trend cannot yet be
disentangled from lr-mismatch (see~\S4.5 and~\S4.8.3).

\section{Future Work}
\label{sec:future-work}
\ \ \ \ \, The following extensions are ordered by estimated effort.

\noindent\textbf{Short-term (within one week).}
(a) Per-$(\text{method},\alpha)$ server-learning-rate sweep at
$\alpha\in\{1.3,1.7\}$ using \path{src/experiments/lr_sweep.py} to
test whether the non-monotonic trend is robust to lr-tuning;
(b) extend the seed count from 3 to 5 for the main $\alpha$ ablation
so that $\pm$ standard deviations tighten below the method-to-method
gap; (c) run the MAC $(k,\gamma)$ sweep at $\alpha\in\{1.1,1.3\}$
to identify an operating region where MAC produces a measurable gain.

\noindent\textbf{Medium-term (within one month).}
Add Rayleigh fading coefficients $h_{n,t}\sim\mathcal{CN}(0,1)$ with
per-client power control $\mathbb{E}|h_{n,t}\Delta_n^t|^2\le\bar P$,
decouple the channel tail index from the optimizer denominator
$\alpha$ so that the two degrees of freedom can be studied
independently, and log per-round mechanism diagnostics
($\|v_t\|$, $\|\tilde g_t\|$, effective step ratio) to produce the
mechanism plot that the current results do not yet include.

\noindent\textbf{Long-term (three to six months).}
Validate the simulator on a software-defined-radio testbed
(USRP X310 and GNU Radio) with multiple transmit clients and a single
server receiver, and replace the one-dimensional closed-form with a
multi-dimensional convergence proof specific to the implemented
fractional-moment update.
